% ---------- Typography ----------
\usepackage{iftex}
\ifluatex
  \usepackage{fontspec}
  \defaultfontfeatures{Ligatures=TeX,Renderer=Harfbuzz}
  % Turon used Minion Pro as both text and display face.  Keep headings,
  % chapter numbers, titles, and body text on the same serif family.
  % Free fallback order is deliberately conservative and book-like:
  % Source Serif 4 -> Libertinus Serif -> EB Garamond -> TeX Gyre Pagella.
  \IfFontExistsTF{Minion Pro}{%
    \setmainfont{Minion Pro}[
      Numbers=OldStyle,
      Ligatures=TeX
    ]%
  }{%
    \IfFontExistsTF{Adobe Source Serif 4}{%
      \setmainfont{Adobe Source Serif 4}[
        Numbers=OldStyle,
        Ligatures=TeX,
        Scale=1.00
      ]%
    }{%
      \IfFontExistsTF{Source Serif 4}{%
        \setmainfont{Source Serif 4}[
          Numbers=OldStyle,
          Ligatures=TeX,
          Scale=1.00
        ]%
      }{%
        \IfFontExistsTF{Libertinus Serif}{%
          \setmainfont{Libertinus Serif}[
            Numbers=OldStyle,
            Ligatures=TeX,
            Scale=1.00
          ]%
        }{%
          \IfFontExistsTF{EB Garamond 12}{%
            \setmainfont{EB Garamond 12}[
              Numbers=OldStyle,
              Ligatures=TeX,
              Scale=0.95
            ]%
          }{%
            \IfFontExistsTF{EB Garamond}{%
              \setmainfont{EB Garamond}[
                Numbers=OldStyle,
                Ligatures=TeX,
                Scale=0.95
              ]%
            }{%
              \IfFontExistsTF{TeX Gyre Pagella}{%
                \setmainfont{TeX Gyre Pagella}[
                  Numbers=OldStyle,
                  Ligatures=TeX,
                  Scale=0.99
                ]%
              }{%
                \setmainfont{lmroman10-regular.otf}[
                  BoldFont=lmroman10-bold.otf,
                  ItalicFont=lmroman10-italic.otf,
                  BoldItalicFont=lmroman10-bolditalic.otf
                ]%
              }%
            }%
          }%
        }%
      }%
    }%
  }

  \IfFontExistsTF{Myriad Pro}{%
    \setsansfont{Myriad Pro}[Ligatures=TeX]
  }{%
    \IfFontExistsTF{Source Sans 3}{%
      \setsansfont{Source Sans 3}[Ligatures=TeX]
    }{%
      \IfFontExistsTF{Source Sans Pro}{%
        \setsansfont{Source Sans Pro}[Ligatures=TeX]
      }{%
        \IfFontExistsTF{TeX Gyre Heros}{%
          \setsansfont{TeX Gyre Heros}[Ligatures=TeX,Scale=.94]
        }{%
          \setsansfont{lmsans10-regular.otf}[
            BoldFont=lmsans10-bold.otf,
            ItalicFont=lmsans10-oblique.otf,
            BoldItalicFont=lmsans10-boldoblique.otf
          ]%
        }%
      }%
    }%
  }

  \IfFontExistsTF{Bera Mono}{%
    \setmonofont{Bera Mono}[Scale=.88]
  }{%
    \IfFontExistsTF{Bitstream Vera Sans Mono}{%
      \setmonofont{Bitstream Vera Sans Mono}[Scale=.88]
    }{%
      \IfFontExistsTF{Bitstream Vera Mono}{%
        \setmonofont{Bitstream Vera Mono}[Scale=.88]
      }{%
        \IfFontExistsTF{DejaVu Sans Mono}{%
          \setmonofont{DejaVu Sans Mono}[Scale=.86]
        }{%
          \IfFontExistsTF{Inconsolata}{%
            \setmonofont{Inconsolata}[Scale=.92]
          }{%
            \setmonofont{lmmono10-regular.otf}[ItalicFont=lmmono10-italic.otf,Scale=.92]%
          }%
        }%
      }%
    }%
  }

  \IfFontExistsTF{DejaVu Sans}{\newfontfamily\turonsymbolfont{DejaVu Sans}[Scale=MatchLowercase]}{%
    \IfFontExistsTF{Noto Sans Symbols 2}{\newfontfamily\turonsymbolfont{Noto Sans Symbols 2}[Scale=MatchLowercase]}{\newfontfamily\turonsymbolfont{Noto Sans}[Scale=MatchLowercase]}}
\else
  \usepackage[T1]{fontenc}
  \usepackage[utf8]{inputenc}
  \usepackage{tgpagella}
  \usepackage[scale=.92]{tgheros}
  \usepackage{beramono}
  \usepackage{pifont}
\fi

% \usepackage[
%   paperwidth=170mm,
%   paperheight=240mm,
%   left=22mm,
%   top=24mm,
%   textwidth=92mm,
%   textheight=190mm,
%   marginparsep=8mm,
%   marginparwidth=33mm,
%   includehead=false,
%   includefoot=false
% ]{geometry}
% ---------- Page and text style ----------
\usepackage[
  a4paper,
  left=25mm,
  right=70mm,
  top=30mm,
  bottom=35mm,
  marginparwidth=34mm,
  marginparsep=9mm,
  heightrounded
]{geometry}
\usepackage{microtype}
\usepackage{booktabs}
\usepackage{graphicx}
\usepackage{import}
\usepackage{xcolor}
\usepackage{enumitem}
\usepackage{amsmath,amssymb}
\usepackage{csquotes}
\usepackage{lettrine}
\usepackage{epigraph}
\usepackage{titlesec}
\usepackage{titletoc}
\usepackage{tocloft}
\usepackage{fancyhdr}
\usepackage{fancyvrb}
\usepackage{caption}
\usepackage{subcaption}
\usepackage{marginnote}
\usepackage{marginfix}
\usepackage{changepage}
\usepackage{etoolbox}
\usepackage{textcase}
\usepackage{pdfpages}
\usepackage{xparse}
\usepackage{xspace}
\usepackage{hyperref}
\usepackage{svg}
\usepackage{tikz}
\usepackage{xparse}
\usetikzlibrary{
  arrows.meta,
  shapes,
  arrows,
  positioning,
  trees,
  shapes.geometric,
  calc,
  fit
}
\usepackage{array}
\providecommand{\Description}[1]{}
\newcolumntype{L}[1]{
  >{\RaggedRight\arraybackslash}p{#1}
}
\usepackage{fontawesome5}
\usepackage{adjustbox}
\usepackage{array}
\usepackage{longtable}
\usepackage{pdflscape}
\usepackage{ragged2e}
\usepackage{tabularx}
\usepackage{float}
\usepackage{placeins}
\newcolumntype{Y}{>{\RaggedRight\arraybackslash}X}
% Front-matter lists can breathe slightly into the fixed right margin.
% This does not mirror page geometry or affect the main text column.
\newenvironment{tocwide}{%
  \begin{adjustwidth}{0pt}{-22mm}%
}{%
  \end{adjustwidth}%
}

% Turon-like accent colours: muted blue for major structural titles;
% restrained wine/magenta for page numbers, sidenote marks, links, and refs.


\definecolor{turonblue}{HTML}{005080}
\definecolor{turonmagenta}{HTML}{8B1E5A}
\definecolor{turongray}{HTML}{555555}
\definecolor{turonlightgray}{HTML}{777777}
% Baseline-based spacing for the chapter-opening grid and heading hierarchy.
% Turon-like refinement comes from a tight baseline rhythm, not large cm-scale gaps.
\newlength{\chapterGapNumTitle}
\newlength{\chapterGapTitleRule}
\newlength{\chapterGapRuleOpening}
\newlength{\sectionBefore}
\newlength{\sectionAfter}
\newlength{\subsectionBefore}
\newlength{\subsectionAfter}
\AtBeginDocument{%
  \setlength{\chapterGapNumTitle}{0.6\baselineskip}%
  \setlength{\chapterGapTitleRule}{0.4\baselineskip}%
  \setlength{\chapterGapRuleOpening}{0.7\baselineskip}%
  \setlength{\sectionBefore}{1.6\baselineskip}%
  \setlength{\sectionAfter}{1.3\baselineskip}%
  \setlength{\subsectionBefore}{1.35\baselineskip}%
  \setlength{\subsectionAfter}{0.95\baselineskip}%
}

% Keep body text true black; reserve grey/colour only for intentional furniture.
\AtBeginDocument{\color{black}}

\hypersetup{
  colorlinks=true,
  linktoc=page,
  linkcolor=turonmagenta,
  citecolor=turonmagenta,
  urlcolor=turonmagenta,
  pdfauthor={\thesisauthor},
  pdftitle={\thesistitle}
}



% Bibliography.
\usepackage[
  backend=biber,
  style=numeric-comp,
  sorting=nyt,
  natbib=true,
  doi=true,
  isbn=false,
  url=true,
  maxcitenames=2,
  maxbibnames=99,
  giveninits=true
]{biblatex}

\addbibresource{references.bib}
\addbibresource{papers/01-scoping-review/scoping-review.bib}
\addbibresource{papers/02-afi/afi.bib}
\addbibresource{papers/03-multisensory-biophilic/multisensory.bib}
\addbibresource{papers/04-bidf/bidf.bib}

% Slightly richer, darker book texture without making the page loose.
\linespread{.96}
\setlength{\parindent}{1em}
\setlength{\parskip}{0pt}
\setlength{\emergencystretch}{1.5em}

% Blank verso/interstitial pages must be visually blank: no running heads,
% footers, or standalone roman numerals. In compact preview mode this behaves
% like \clearpage; in final dissertation mode it enforces recto openings while
% keeping the fixed, non-mirrored Tufte text column.
\makeatletter
\@ifundefined{iffinaldissertation}{\newif\iffinaldissertation\finaldissertationfalse}{}
\newif\ifturonstartedpages
\AddToHook{shipout/after}{\global\turonstartedpagestrue}
\renewcommand{\cleardoublepage}{%
  \clearpage
  \iffinaldissertation
    % In oneside final mode, keep recto openings without changing geometry.
    % The guard prevents a blank page before the title page at document start.
    \ifturonstartedpages
      \ifodd\c@page
        \hbox{}%
        \thispagestyle{empty}%
        \newpage
        \if@twocolumn\hbox{}\newpage\fi
      \fi
    \fi
  \else
    \if@twoside
      \ifodd\c@page
      \else
        \hbox{}%
        \thispagestyle{empty}%
        \newpage
        \if@twocolumn\hbox{}\newpage\fi
      \fi
    \fi
  \fi}
\makeatother

% ---------- Heading hierarchy ----------

% Use lining figures for large chapter numbers even when body text uses old-style figures.
\newcommand{\turonchapnum}[1]{{\ifluatex\addfontfeatures{Numbers=Lining}\fi\bfseries #1}}
% Chapter openings: large black number only, no word "Chapter".
% Baseline grid: number--title = .6 baseline; title--rule = .4 baseline;
% rule--opening paragraph = .7 baseline. The rule sits close under the title.
\titleformat{\chapter}[display]
  {\normalfont\raggedright}
  {\color{black}\fontsize{76}{78}\selectfont\turonchapnum{\thechapter}}
  {\chapterGapNumTitle}
  {\normalfont\Large\itshape\color{black!76}}
  [\vspace{\chapterGapTitleRule}\noindent\textcolor{black!24}{\rule{\textwidth}{0.3pt}}\vspace{\chapterGapRuleOpening}]
\titlespacing*{\chapter}{0pt}{-22pt}{0pt}

% Unnumbered front-matter chapters: quiet italic heading.
\titleformat{name=\chapter,numberless}[display]
  {\normalfont\raggedright}
  {}
  {0pt}
  {\normalfont\Large\itshape\color{black!86}}
\titlespacing*{name=\chapter,numberless}{0pt}{1.3em}{1.6em}


\newcommand{\designbox}[3]{%
\begin{figure}[!htbp]
\centering
\fbox{%
\begin{minipage}{0.95\linewidth}
\textbf{#1}

\vspace{4pt}
#2
\end{minipage}}
\Description{#3}
\end{figure}
}




% ---------- Publication list ----------
\newenvironment{publicationlist}{%
  \begingroup
  \fontsize{9.3}{10.8}\selectfont
  \RaggedRight
  \setlength{\parindent}{0pt}
  \setlength{\parskip}{.48\baselineskip}
}{%
  \par
  \endgroup
}

\newcommand{\pubentry}[1]{%
  \par\noindent
  \hangindent=1.15em
  \hangafter=1
  #1\par
}

\newcommand{\pubaward}[1]{%
  \nobreak\hspace{.28em}%
  {\scriptsize\textcolor{turonblue}{\faAward}}%
  \hspace{.15em}%
  {\scriptsize\itshape #1}%
}

% ---------- Section heading style ----------
\setcounter{secnumdepth}{2}

\newcommand{\turonsectiontitle}[1]{%
  \textls[22]{\MakeTextUppercase{#1}}%
}

\newcommand{\turonsubsectiontitle}[1]{%
  \textit{#1}%
}

% Sections: body-size, uppercase small caps.
\titleformat{\section}
  {\normalfont\normalsize\scshape\raggedright\color{black}}
  {\normalfont\small\upshape\thesection}
  {1.0em}
  {\turonsectiontitle}

\titlespacing*{\section}
  {0pt}
  {1.85\baselineskip}
  {1.05\baselineskip}

% Subsections: numbered, italic, quiet.
% This is the deepest formal heading level in the dissertation body.
\titleformat{\subsection}
  {\normalfont\normalsize\itshape\raggedright\color{black}}
  {\normalfont\small\upshape\thesubsection}
  {.85em}
  {\turonsubsectiontitle}

\titlespacing*{\subsection}
  {0pt}
  {1.45\baselineskip}
  {.75\baselineskip}

% Subsubsections: available, but unnumbered and rarely used.
% Prefer \paragraph for local labelled blocks.
\titleformat{\subsubsection}[hang]
  {\normalfont\normalsize\itshape\raggedright\color{black}}
  {}
  {0pt}
  {}

\titlespacing*{\subsubsection}
  {0pt}
  {1.20\baselineskip}
  {.55\baselineskip}

% Paragraph-level minor heading with Turon-like grey arrow.
% Use for steps, themes, practices, opportunities, values, and local analytic moves.

% ---------- Turon-like flat arrow ----------
\newcommand{\turonarrow}{%
  \tikz[baseline=-0.48ex]{%
    \fill[black!35]
      (0,0) --
      (0.95ex,0.50ex) --
      (0,1.00ex) --
      cycle;
  }%
}
\titleformat{\paragraph}[hang]
  {\normalfont\normalsize\scshape\raggedright\color{black}}
  {}
  {0pt}
  {\turonarrow\hspace{.65em}}

\titlespacing*{\paragraph}
  {0pt}
  {1.05\baselineskip}
  {.45\baselineskip}

\titlespacing*{\paragraph}
  {0pt}
  {1.05\baselineskip}
  {.45\baselineskip}
  
% Part pages: quiet "Part I" plus blue small-caps/all-caps title, centered mid-page.
\newcommand{\turonparttitle}[1]{\textls[55]{\MakeTextUppercase{#1}}}
\titleformat{\part}[display]
  {\thispagestyle{empty}\normalfont\centering}
  {\normalfont\small Part \Roman{part}}
  {1.8em}
  {\normalfont\Large\scshape\color{turonblue}\turonparttitle}
\titlespacing*{\part}{0pt}{0.30\textheight}{0.38\textheight}

% ---------- Front matter helper ----------
\newcommand{\frontmatterheading}[1]{%
  \chapter*{#1}%
  \markboth{\MakeTextLowercase{#1}}{\MakeTextLowercase{#1}}%
}

% Paragraph opening marker, like Turon's compact triangular synopsis lead-in.
% Use an explicit Unicode glyph with a known symbol font under LuaLaTeX so it cannot
% degrade to an I, dot, or math fallback.
\ifluatex
  \DeclareRobustCommand{\synopsismarker}{{\turonsymbolfont\symbol{"25B8}}}
\else
  \DeclareRobustCommand{\synopsismarker}{\ding{228}}
\fi
\DeclareRobustCommand{\openingmark}{%
  \noindent\llap{\raisebox{.01em}{\scriptsize\color{turonblue!55!black}\synopsismarker}\hspace{.7em}}\ignorespaces%
}

% A compact opening synopsis: marker + paragraph flow, not a labelled box.
% It always begins below the chapter epigraph, rather than visually merging with it.
\newenvironment{synopsis}{%
  \par\addvspace{0pt}%
  \begingroup
  \fontsize{10.5}{12.4}\selectfont
  \color{black}%
  \openingmark\ignorespaces
}{%
  \par\endgroup\addvspace{0pt}%
}


% ---------- Turon-style small-caps lead-ins ----------
% The inspiration uses small caps, not bold all-caps.
% Keep the source text lowercase/titlecase; this macro renders it as small caps.
\newcommand{\turonlead}[1]{%
  {\normalfont\scshape\textls[55]{\MakeTextLowercase{#1}}}%
}

\newcommand{\openingpara}[2]{%
  \par\addvspace{0pt}%
  \noindent\llap{\raisebox{.01em}{\scriptsize\color{turonblue!55!black}\synopsismarker}\hspace{.7em}}%
  \turonlead{#1} #2\par\addvspace{0pt}%
}

\newcommand{\synopsislead}[2]{%
  \turonlead{#1} #2%
}
% Uppercase/small-caps lead phrase for use inside the synopsis environment.
% Use when the synopsis environment already provides the blue opening marker.
% \newcommand{\synopsislead}[2]{%
%   {\scshape\textls[8]{\MakeTextUppercase{#1}}} #2%
% }

% ----------chapter closing page ----------
% A quiet dissertation-level coda placed after each paper chapter.
% The optional argument is retained for compatibility but is not printed.

\newcommand{\chapterclosinglead}[2]{%
  \par\addvspace{0pt}%
  \noindent\turonlead{#1} #2\par
}

\NewDocumentEnvironment{chapterclosing}{O{} m}
  {%
    \clearpage
    \thispagestyle{plain}%
    \begingroup
    \vspace*{0.18\textheight}%

    {\noindent\normalfont\Large\itshape\color{black!76}#2\par}%
    \vspace{.45\baselineskip}%

    \noindent\textcolor{black!24}{\rule{\textwidth}{0.3pt}}%
    \par\vspace{.9\baselineskip}%

    \fontsize{10.5}{12.4}\selectfont
    \RaggedRight
    \setlength{\parindent}{1em}%
    \setlength{\parskip}{0pt}%
  }
  {%
    \par
    \endgroup
    \clearpage
  }

% Optional final bridge paragraph, now plain rather than labelled.
\newcommand{\chaptercarryforward}[1]{%
  \par\addvspace{.75\baselineskip}%
  {\footnotesize\RaggedRight #1\par}%
}

% ---------- Chapter source note ----------
% Used at the end of chapter-opening pages to distinguish the publication
% basis of an included paper chapter from the main synopsis text.
\newcommand{\chaptersource}[1]{%
  \par\addvspace{1.15\baselineskip}%
  \noindent\textcolor{black!24}{\rule{.52\linewidth}{0.3pt}}%
  \par\vspace{.45\baselineskip}%
  {\footnotesize\RaggedRight
    {\normalfont\scshape\textls[35]{\textcolor{turonblue}{Chapter source}}}\par
    \vspace{.22\baselineskip}%
    #1\par
  }%
  \par\addvspace{.25\baselineskip}%
}


% ---------- Sidenotes and margin material ----------
% Keep all notes in the fixed right margin. Do not mirror them on even pages.
\setlength{\marginparpush}{6pt}
\setlength{\marginparsep}{8.5mm}
\setlength{\marginparwidth}{34mm}
\renewcommand*{\marginfont}{%
  \fontsize{6.2}{7.25}\selectfont
  \RaggedRight
  \setlength{\emergencystretch}{1em}%
}
\newcommand*{\epigraphmarginfont}{\fontsize{7.2}{8.6}\selectfont\raggedleft}
\renewcommand*{\raggedleftmarginnote}{\raggedright}
\renewcommand*{\raggedrightmarginnote}{\raggedright}
\newcounter{sidenote}
\newcommand{\turonnotemark}{\textsuperscript{\tiny\textcolor{turonmagenta!65!black}{\thesidenote}}}
\newcommand{\turonmarginmark}{\textsuperscript{\tiny\textcolor{turonmagenta!65!black}{\thesidenote}}}
\newcommand{\turonmarginnote}[1]{%
  \normalmarginpar%
  \marginpar{%
    \marginfont
    \setlength{\parindent}{0pt}%
    \setlength{\parskip}{0pt}%
    \RaggedRight
    \turonmarginmark\,\textcolor{black}{#1}%
  }%
}
\newcommand{\sidenote}[1]{%
  \refstepcounter{sidenote}%
  \turonnotemark%
  \turonmarginnote{#1}%
}
\newcommand{\sidenotedown}[2]{%
  \refstepcounter{sidenote}%
  \turonnotemark%
  \marginnote{\marginfont\turonmarginmark\,\textcolor{black}{#2}}[#1]%
}
\newcommand{\marginterm}[2]{\textbf{#1}\marginpar{\marginfont\textbf{#1.} #2}}
% Chapter epigraphs sit in the fixed right margin, aligned with the opening
% paragraph/rule area rather than occupying the main text flow. Using marginpar
% keeps subsequent sidenotes from colliding with the epigraph.
\newcommand{\sidequote}[2]{%
  \par\addvspace{0pt}%
  \normalmarginpar%
  \marginpar{\parbox[t]{\marginparwidth}{\epigraphmarginfont\setlength{\parindent}{0pt}\itshape ``#1''\par\vspace{.35em}{\normalfont\scriptsize\raggedleft---#2\par}}}%
  \par\addvspace{0pt}%
}
\newcommand{\chapterepigraph}[2]{\sidequote{#1}{#2}}
\newcommand{\sidenoteup}[2]{%
  \refstepcounter{sidenote}%
  \turonnotemark%
  \marginnote{%
    \marginfont
    \turonmarginmark\,\textcolor{black}{#2}%
  }[-#1]%
}
\newsavebox{\marginfigurebox}
\newenvironment{marginfigure}{%
  \begin{lrbox}{\marginfigurebox}\begin{minipage}{\marginparwidth}\marginfont\centering
}{%
  \end{minipage}\end{lrbox}\marginnote{\usebox{\marginfigurebox}}[0pt]%
}


% ---------- Figure and table width system ----------
% Wide material spans the main text column + fixed right-margin note column.
% It does not mirror and does not change the page geometry.
\newlength{\turonwideoffset}
\AtBeginDocument{%
  \setlength{\turonwideoffset}{\dimexpr\marginparsep+\marginparwidth\relax}%
}

\newenvironment{fullwidth}{%
  \begin{adjustwidth}{0pt}{-\turonwideoffset}%
}{%
  \end{adjustwidth}%
}


% Main-column figure: use for ordinary, simple figures.
\NewDocumentEnvironment{mainfigure}{O{tbp}}
  {%
    \figure[#1]%
    \centering
    \captionsetup{width=\linewidth}%
  }
  {%
    \endfigure
  }

% Wide figure: use for dense diagrams, frameworks, Sankeys, system diagrams.
% Spans main text + right-margin note column.
\NewDocumentEnvironment{widefigure}{O{tbp}}
  {%
    \figure[#1]%
    \fullwidth
    \centering
    \captionsetup{width=\linewidth}%
  }
  {%
    \endfullwidth
    \endfigure
  }

% Main-column table: use for compact summary tables.
\NewDocumentEnvironment{maintable}{O{tbp}}
  {%
    \table[#1]%
    \centering
    \small
    \captionsetup{width=\linewidth}%
  }
  {%
    \endtable
  }

% Wide table: use for tables that need the text column + margin column.
\NewDocumentEnvironment{widetable}{O{tbp}}
  {%
    \table[#1]%
    \fullwidth
    \centering
    \small
    \captionsetup{width=\linewidth}%
  }
  {%
    \endfullwidth
    \endtable
  }

% ---------- Float placement tuning ----------
% Allow figures and tables to appear nearer to where they are introduced.
\renewcommand{\topfraction}{0.90}
\renewcommand{\bottomfraction}{0.80}
\renewcommand{\textfraction}{0.07}
\renewcommand{\floatpagefraction}{0.75}
\renewcommand{\dbltopfraction}{0.90}
\renewcommand{\dblfloatpagefraction}{0.75}

\setcounter{topnumber}{4}
\setcounter{bottomnumber}{3}
\setcounter{totalnumber}{6}

% ---------- Appendix table helpers ----------
% Use for dense appendix tables. Avoid resizebox unless absolutely necessary.
\newcommand{\appendixtablefont}{\fontsize{8.6}{10.2}\selectfont}
\newcommand{\appendixwidefont}{\fontsize{9.2}{11}\selectfont}

\newcommand{\tighttable}{%
  \setlength{\tabcolsep}{3.5pt}%
  \renewcommand{\arraystretch}{1.08}%
}

\newcommand{\verytighttable}{%
  \setlength{\tabcolsep}{2.8pt}%
  \renewcommand{\arraystretch}{1.04}%
}
% Lists: compact, solid bullet, no report-like gaps.
\setlist[itemize]{leftmargin=1.2em,itemsep=.15em,topsep=.35em,parsep=0pt,partopsep=0pt,label={\color{turonblue}$\bullet$}}
\setlist[enumerate]{leftmargin=1.45em,itemsep=.15em,topsep=.35em,parsep=0pt,partopsep=0pt}


% ---------- Captions ----------
% Turon/Tufte-like captions: a small left label column and prose caption to the right.
% This keeps long published captions readable without shortening them.

\newlength{\turoncaptionlabelwidth}
\setlength{\turoncaptionlabelwidth}{20 mm}

\DeclareCaptionLabelFormat{turonlabel}{%
  \textcolor{turonblue}{\scshape #1~#2.}%
}

\DeclareCaptionFormat{turoncaptionformat}{%
  \noindent
  \makebox[\turoncaptionlabelwidth][l]{#1}%
  \begin{minipage}[t]{\dimexpr\linewidth-\turoncaptionlabelwidth\relax}
    \RaggedRight
    #3%
  \end{minipage}%
  \par
}

\captionsetup{
  format=turoncaptionformat,
  labelformat=turonlabel,
  labelsep=none,
  justification=RaggedRight,
  singlelinecheck=false,
  font={scriptsize},
  labelfont={normalfont},
  textfont={normalfont},
  margin=0pt,
  indention=0pt,
  aboveskip=.55\baselineskip,
  belowskip=.55\baselineskip
}

\captionsetup[figure]{name=Figure}
\captionsetup[table]{name=Table}


% Code blocks.
\fvset{fontsize=\small,frame=single,rulecolor=\color{black!15},framesep=3mm}

% ---------- Contents and lists ----------
\setcounter{tocdepth}{1}
\renewcommand{\contentsname}{Contents}
\renewcommand{\listfigurename}{Figures}
\renewcommand{\listtablename}{Tables}

\renewcommand{\cfttoctitlefont}{\normalfont\Large\itshape\color{black!86}}
\renewcommand{\cftloftitlefont}{\normalfont\Large\itshape\color{black!86}}
\renewcommand{\cftlottitlefont}{\normalfont\Large\itshape\color{black!86}}
\renewcommand{\cftaftertoctitle}{\par\vspace{-.25em}\noindent\rule{\textwidth}{.35pt}\vspace{1.0em}}
\renewcommand{\cftafterloftitle}{\par\vspace{-.25em}\noindent\rule{\textwidth}{.35pt}\vspace{1.0em}}
\renewcommand{\cftafterlottitle}{\par\vspace{-.25em}\noindent\rule{\textwidth}{.35pt}\vspace{1.0em}}
\setlength{\cftbeforetoctitleskip}{0pt}
\setlength{\cftbeforeloftitleskip}{0pt}
\setlength{\cftbeforelottitleskip}{0pt}
\setlength{\cftaftertoctitleskip}{.55em}
\setlength{\cftafterloftitleskip}{.5em}
\setlength{\cftafterlottitleskip}{.5em}

% TOC entries: compact, Turon/classicthesis-like hierarchy.
\newcommand{\tocnohyphen}{%
  \hyphenpenalty=10000%
  \exhyphenpenalty=10000%
  \pretolerance=10000%
  \tolerance=2200%
  \emergencystretch=4.0em%
  \relax
}
% Parts: blue, spaced, small caps.
\renewcommand{\cftpartfont}{%
  \tocnohyphen\normalfont\small\scshape\bfseries\lsstyle\color{turonblue}%
}
\renewcommand{\cftpartpagefont}{%
  \normalfont\small\bfseries\color{turonmagenta}%
}

% Chapters: tracked small caps, clear but not oversized.
\renewcommand{\cftchapfont}{%
  \tocnohyphen\normalfont\small\scshape\bfseries\lsstyle\color{black}%
}
\renewcommand{\cftchappagefont}{%
  \normalfont\small\bfseries\color{turonmagenta}%
}

% Sections: quieter than chapters.
\renewcommand{\cftsecfont}{%
  \tocnohyphen\normalfont\small\color{black}%
}
\renewcommand{\cftsecpagefont}{%
  \normalfont\small\color{turonmagenta}%
}

% Subsections: available if tocdepth is later changed to 2.
\renewcommand{\cftsubsecfont}{%
  \tocnohyphen\normalfont\footnotesize\color{black}%
}
\renewcommand{\cftsubsecpagefont}{%
  \normalfont\footnotesize\color{turonmagenta}%
}

% Slightly calmer dot leaders.
\renewcommand{\cftdotsep}{4.8}
\renewcommand{\cftchapleader}{\cftdotfill{\cftdotsep}}
\renewcommand{\cftsecleader}{\cftdotfill{\cftdotsep}}
\renewcommand{\cftsubsecleader}{\cftdotfill{\cftdotsep}}

% More air before major units.
\setlength{\cftbeforepartskip}{1.15em}
\setlength{\cftbeforechapskip}{.46em}
\setlength{\cftbeforesecskip}{.09em}
\setlength{\cftbeforesubsecskip}{.02em}

% Indentation.
\setlength{\cftpartindent}{0pt}
\setlength{\cftchapindent}{0pt}
\setlength{\cftsecindent}{1.95em}
\setlength{\cftsubsecindent}{3.6em}

% Number widths.
\setlength{\cftpartnumwidth}{1.9em}
\setlength{\cftchapnumwidth}{2.0em}
\setlength{\cftsecnumwidth}{3.0em}
\setlength{\cftsubsecnumwidth}{3.7em}

% Figure/table lists: Turon-like prefix and magenta page numbers.
\renewcommand{\cftfigpresnum}{Figure~}
\renewcommand{\cfttabpresnum}{Table~}
\setlength{\cftfignumwidth}{4.8em}
\setlength{\cfttabnumwidth}{4.4em}
\renewcommand{\cftfigfont}{\normalfont\small}
\renewcommand{\cftfigpagefont}{\normalfont\small\color{turonmagenta}}
\renewcommand{\cfttabfont}{\normalfont\small}
\renewcommand{\cfttabpagefont}{\normalfont\small\color{turonmagenta}}
\renewcommand{\cftfigleader}{\cftdotfill{\cftdotsep}}
\renewcommand{\cfttableader}{\cftdotfill{\cftdotsep}}
\setlength{\cftbeforefigskip}{.12em}
\setlength{\cftbeforetabskip}{.12em}

% ---------- Stable Tufte-style running heads ----------
% The header does not mirror across odd/even pages: text column stays fixed.
% Page numbers sit quietly in the bottom-right footer; magenta is reserved for links/TOC.
\setlength{\headheight}{17pt}
\pagestyle{fancy}
\fancyhf{}
\fancyfootoffset[R]{\dimexpr\marginparsep+\marginparwidth\relax}
\renewcommand{\chaptermark}[1]{\markboth{\thechapter\quad\MakeTextLowercase{#1}}{}}
\renewcommand{\sectionmark}[1]{\markright{\thesection\quad\MakeTextLowercase{#1}}}
\fancyhead[L]{\scriptsize\scshape\color{turongray}\nouppercase{\leftmark}}
\fancyfoot[R]{\footnotesize\thepage}
\renewcommand{\headrulewidth}{0pt}
\fancypagestyle{plain}{\fancyhf{}\fancyfoot[R]{\footnotesize\thepage}\renewcommand{\headrulewidth}{0pt}}
\fancypagestyle{empty}{\fancyhf{}\renewcommand{\headrulewidth}{0pt}}

% Stable reference helpers. Each macro prints the word and number as one link
% with a non-breaking space inside the phrase, then lets \xspace preserve
% the following word space in constructions such as "\chapref{ch:a} and".
% Prefer the explicit space in source, but these also survive tight usages.
\DeclareRobustCommand{\chapref}[1]{\hyperref[#1]{Chapter~\ref*{#1}}\xspace}
\DeclareRobustCommand{\secref}[1]{\hyperref[#1]{Section~\ref*{#1}}\xspace}
\DeclareRobustCommand{\figref}[1]{\hyperref[#1]{Figure~\ref*{#1}}\xspace}
\DeclareRobustCommand{\tabref}[1]{\hyperref[#1]{Table~\ref*{#1}}\xspace}

% Useful placeholders.
\newcommand{\todo}[1]{\marginnote{\marginfont\color{turonmagenta}TODO: #1}[0pt]}


% ---------- Research questions and research contributions ----------
% Use these to create stable, referable labels such as RQ1, RQ2, RC1, RC2.

\newcounter{researchquestion}
\newcounter{researchcontribution}

\newcommand{\rqtag}[1]{%
  {\normalfont\normalsize\textcolor{turonblue}{RQ~#1}}%
}

\newcommand{\rctag}[1]{%
  {\normalfont\normalsize\textcolor{turonblue}{RC~#1}}%
}

\newenvironment{researchquestions}
  {%
    \setcounter{researchquestion}{0}%
    \begin{list}{}{%
      \setlength{\leftmargin}{4.7em}%
      \setlength{\labelwidth}{3.7em}%
      \setlength{\labelsep}{1em}%
      \setlength{\itemsep}{0.65\baselineskip}%
      \setlength{\parsep}{0pt}%
      \setlength{\topsep}{0.55\baselineskip}%
    }%
  }
  {\end{list}}

\newcommand{\researchquestion}[2]{%
  \refstepcounter{researchquestion}%
  \item[\rqtag{\theresearchquestion}]%
  \label{rq:#1}%
  #2%
}

\newcommand{\rqref}[1]{%
  \hyperref[rq:#1]{\rqtag{\ref*{rq:#1}}}%
}

\newenvironment{researchcontributions}
  {%
    \setcounter{researchcontribution}{0}%
    \begin{list}{}{%
      \setlength{\leftmargin}{4.7em}%
      \setlength{\labelwidth}{3.7em}%
      \setlength{\labelsep}{1em}%
      \setlength{\itemsep}{0.75\baselineskip}%
      \setlength{\parsep}{0pt}%
      \setlength{\topsep}{0.55\baselineskip}%
    }%
  }
  {\end{list}}

\newcommand{\researchcontribution}[2]{%
  \refstepcounter{researchcontribution}%
  \item[\rctag{\theresearchcontribution}]%
  \label{rc:#1}%
  #2%
}

\newcommand{\rcref}[1]{%
  \hyperref[rc:#1]{\rctag{\ref*{rc:#1}}}%
}

\newcommand{\orqref}{%
  \hyperref[rq:overarching]{\normalfont\scshape\bfseries\textcolor{turonblue}{ORQ}}%
}